% -------------------------------------------------------------------------
% WBS ID: RES-VER-TST
% Deliverable: Component, Conservation and Interface Test Design
% Status: Not started
% Evidence status: Not assessed
% Review status: Not reviewed
% Last updated:
% Dependencies:
% Downstream interfaces:
% -------------------------------------------------------------------------

\section{RES-VER-TST --- Component, Conservation and Interface Test Design}
\label{sec:res-ver-tst}

\subsection{Purpose and Scope}
\label{sec:res-ver-tst-purpose}

% TODO: Define what this Deliverable covers and explicitly excludes.

\subsection{Research Questions}
\label{sec:res-ver-tst-research-questions}

% TODO: State the questions that must be answered before the Deliverable
% can be accepted.

\subsection{Terminology and Definitions}
\label{sec:res-ver-tst-definitions}

% TODO: Define the technical terminology, symbols, quantities and
% classifications used in this Deliverable.

\subsection{Literature Evidence}
\label{sec:res-ver-tst-literature-evidence}

% TODO: Record evidence from peer-reviewed papers, authoritative datasets,
% standards, technical reports and primary documentation.
%
% Do not create a detached literature-summary list. Explain how each source
% supports, contradicts or limits a project decision.

\subsubsection{Qin et al. (2018)}
\label{sec:res-ver-tst-evidence-qin-2018}

\textcite{qin2018comparison} demonstrated that a local
three-dimensional model should be assessed against depth-dependent
velocity and direct structural-force measurements rather than only
against free-surface elevation.

\subsubsection{Watanabe et al. (2022)}
\label{sec:res-ver-tst-evidence-watanabe-2022}

\textcite{watanabe2022validation} compared multiple numerical
predictions against withheld laboratory measurements. The blind
comparison showed that three-dimensional models could reproduce
nonhydrostatic and short-wave behaviour but also produced substantial
variability in velocity and force. Smaller cells did not systematically
reduce this variability.

\subsection{Local Tsunami-Impact Test Sequence}
\label{sec:res-ver-tst-local-impact-sequence}

The selected parent equations shall be validated as part of a complete
modelling workflow. Agreement in water level shall not be treated as
evidence of accurate pressure, force, or structural loading.

The local model shall be assessed progressively using:

\begin{enumerate}
  \item incident-wave generation without a structure;
  \item free propagation and transformation within the local domain;
  \item impact on an isolated impermeable structure;
  \item pressure and force on a simple canonical geometry;
  \item overtopping and run-up on a barrier;
  \item separation, wake, and downstream-flow validation;
  \item interaction with multiple or finite-width structures;
  \item one-way structural load transfer;
  \item two-way FSI where activated.
\end{enumerate}

Each stage shall isolate a limited set of physical and numerical
mechanisms before additional complexity is introduced.

\subsection{One-Way and Simultaneously Coupled Model Comparison}
\label{sec:res-ver-tst-coupling-mode-comparison}

\subsubsection{Test Objective}
\label{sec:res-ver-tst-coupling-mode-objective}

The test shall determine whether neglecting feedback from the local
three-dimensional barrier domain produces a material change in the
quantities used to assess candidate barrier performance. The comparison
shall distinguish:

\begin{enumerate}
  \item the one-way replay model, in which a saved two-dimensional
    state drives the local three-dimensional model without reverse
    feedback;
  \item the simultaneously coupled model, in which the retained
    two-dimensional buffer and local three-dimensional model
    exchange information during the impact calculation.
\end{enumerate}

\subsubsection{Test Configuration}
\label{sec:res-ver-tst-coupling-mode-configuration}

Both branches shall begin from the same immutable pre-impact checkpoint
and shall use identical:

\begin{itemize}
  \item environmental scenario;
  \item barrier geometry;
  \item material and boundary assumptions;
  \item interface location;
  \item two-dimensional and three-dimensional meshes;
  \item physical simulation interval;
  \item output locations;
  \item numerical tolerances;
  \item quantity definitions.
\end{itemize}

The comparison shall be repeated for at least:

\begin{itemize}
  \item the reference barrier;
  \item a weakly reflecting candidate;
  \item a strongly reflecting or obstructive candidate;
  \item more than one interface location.
\end{itemize}

\subsubsection{Concurrent Execution}
\label{sec:res-ver-tst-coupling-mode-concurrency}

The two branches may be launched concurrently to reduce elapsed
comparison time. Each branch shall receive a separate copy of the saved
checkpoint and separate mutable simulation state. A synchronisation
barrier shall prevent the comparison and acceptance logic from
executing until both simulations have completed successfully. A failed
branch shall invalidate the comparison rather than permitting the
surviving result to be accepted in isolation. The assigned cores,
threads, accelerators, and memory shall be recorded. Concurrent
execution shall not oversubscribe the available hardware or change the
numerical configuration relative to separately executed reference runs.

\subsubsection{Measured Quantities}
\label{sec:res-ver-tst-coupling-mode-quantities}

The comparison shall include:

\begin{itemize}
  \item free-surface elevation upstream and downstream of the barrier;
  \item reflected-wave amplitude and energy;
  \item transmitted wave energy;
  \item normal mass flux across the interface;
  \item momentum flux across the interface;
  \item peak barrier force;
  \item barrier-force impulse;
  \item peak pressure;
  \item overtopping volume;
  \item downstream maximum depth and velocity;
  \item regional run-up or inundation where affected;
  \item wall-clock time, processor time, and memory consumption.
\end{itemize}

\subsubsection{Error Metrics}
\label{sec:res-ver-tst-coupling-mode-error-metrics}

For a scalar quantity of interest $q$, the relative difference shall be
evaluated using:

\begin{align}
  \delta_q
  &=
  \frac{
    \left| q_{\mathrm{coupled}} - q_{\mathrm{one\mbox{-}way}} \right|
  }{
    \max \left( \left|q_{\mathrm{coupled}}\right|, q_{\mathrm{scale}} \right)
  }
  \label{eq:res-ver-tst-coupling-relative-difference}
\end{align}

where $q_{\mathrm{scale}}$ is a non-zero reference scale used to avoid
unstable normalisation for quantities close to zero. For a
time-dependent quantity $q(t)$, the normalised difference may be
evaluated using:

\begin{align}
  \delta_{q,2}
  &=
  \frac{
    \left\| q_{\mathrm{coupled}} - q_{\mathrm{one\mbox{-}way}} \right\|_{L^2}
  }{
    \max \left( \left\| q_{\mathrm{coupled}} \right\|_{L^2}, q_{\mathrm{scale}} \right)
  }
  \label{eq:res-ver-tst-coupling-time-series-difference}
\end{align}

Peak-value, phase, impulse, and spatial-location errors shall be
reported separately rather than being concealed within one aggregate
norm.

\subsubsection{Tolerance and Acceptance Criteria}
\label{sec:res-ver-tst-coupling-mode-acceptance}

The one-way replay model may be accepted for broad candidate screening
when:

\begin{itemize}
  \item all safety-critical and optimisation quantities remain within
    their prescribed fidelity tolerances;
  \item interface mass and momentum imbalances remain acceptable;
  \item reflected-wave differences do not materially alter the upstream
    or downstream response;
  \item the conclusion remains stable under movement of the interface;
  \item the ranking of candidate geometries is unchanged.
\end{itemize}

The simultaneously coupled model shall be required when any
safety-critical quantity exceeds its tolerance, when candidate rankings
change, or when barrier-generated feedback materially modifies the
regional response. The acceptance tolerances shall be defined after the
characteristic magnitudes, numerical errors, and decision sensitivity of
each quantity have been established. They shall not be selected solely
to favour the cheaper model.

\subsubsection{Execution Handoff}
\label{sec:res-ver-tst-coupling-mode-handoff}

The Software workstream shall automate:

\begin{enumerate}
  \item restoration of the common checkpoint;
  \item creation of isolated branch states;
  \item concurrent or sequential execution;
  \item branch-completion synchronisation;
  \item failure detection;
  \item metric extraction;
  \item comparison against the configured tolerances;
  \item generation of a reproducible test report.
\end{enumerate}

The Optimisation workstream shall use the approved fidelity policy to
select the solver mode for candidate screening and final confirmation.

\subsection{Interface-Location Sensitivity Test}
\label{sec:res-ver-tst-interface-location}

\subsubsection{Test Objective}
\label{sec:res-ver-tst-interface-location-objective}

The test shall determine whether the selected
two-dimensional--three-dimensional interfaces lie within regions where
the assumptions of their receiving models are valid.

Separate studies shall be performed for:

\begin{enumerate}
  \item the upstream two-dimensional-to-three-dimensional lifting
    interface;
  \item the downstream three-dimensional-to-two-dimensional
    restriction interface;
  \item the bidirectional interface used by the simultaneously coupled
    configuration.
\end{enumerate}

\subsubsection{Test Configuration}
\label{sec:res-ver-tst-interface-location-configuration}

At least three candidate interface positions shall be examined:

\begin{enumerate}
  \item a position close to the local barrier-interaction region;
  \item a position near the expected transition in model validity;
  \item a position farther into the region where the receiving model
    is expected to be valid.
\end{enumerate}

The candidate locations shall be selected with reference to:

\begin{itemize}
  \item bathymetric and topographic gradients;
  \item geometric contractions and expansions;
  \item predicted breaking and overtopping regions;
  \item vertical velocity and shear;
  \item free-surface fragmentation;
  \item recirculation and wake length;
  \item reflected-wave propagation;
  \item wetting--drying behaviour;
  \item proximity to artificial boundaries.
\end{itemize}

All other physical and numerical inputs shall remain unchanged between
the interface-location cases.

\subsubsection{Measured Quantities}
\label{sec:res-ver-tst-interface-location-quantities}

The test shall compare:

\begin{itemize}
  \item free-surface time histories;
  \item depth-integrated discharge;
  \item incident, reflected, and transmitted wave amplitudes;
  \item incident, reflected, and transmitted energy;
  \item peak barrier force;
  \item barrier-force impulse;
  \item peak pressure and pressure impulse;
  \item overtopping volume;
  \item downstream depth and velocity;
  \item run-up and inundation extent;
  \item interface mass imbalance;
  \item interface momentum imbalance;
  \item amplitude and phase spectra.
\end{itemize}

\subsubsection{Error Metrics}
\label{sec:res-ver-tst-interface-location-error}

For a quantity of interest $q$ evaluated using neighbouring interface
positions $\Gamma_i$ and $\Gamma_{i+1}$, the relative location
sensitivity shall be defined as:

\begin{align}
  \delta_{\Gamma,q}^{(i)}
  &=
  \frac{
    \left|
    q_{\Gamma_{i+1}}
    -
    q_{\Gamma_i}
    \right|
  }{
    \max
    \left(
      \left|q_{\Gamma_{i+1}}\right|,
      q_{\mathrm{scale}}
    \right)
  }
  \label{eq:res-ver-tst-interface-location-sensitivity}
\end{align}

where $q_{\mathrm{scale}}$ is a non-zero characteristic scale.

For a transferred waveform $\eta(t)$, the amplitude spectrum shall be
evaluated using:

\begin{align}
  \widehat{\eta}(f)
  &=
  \mathcal{F}
  \left\{
    \eta(t)
  \right\}
  \label{eq:res-ver-tst-interface-spectrum}
\end{align}

Spectral amplitude and phase differences shall be reported alongside
time-domain norms.

\subsubsection{Tolerance and Acceptance Criteria}
\label{sec:res-ver-tst-interface-location-acceptance}

An interface position may be accepted when:

\begin{itemize}
  \item moving the interface farther into the receiving model's valid
    region produces no material change in the engineering
    quantities of interest;
  \item mass and momentum transfer errors remain within their
    prescribed tolerances;
  \item the transferred waveform retains the required amplitude,
    phase, and frequency content;
  \item candidate-design rankings remain unchanged;
  \item no material artificial-boundary contamination is detected;
  \item the receiving model reproduces the donor behaviour over an
    overlapping comparison region.
\end{itemize}

The cheapest acceptable interface position shall be selected. A
location shall not be chosen solely to minimise the size of the
three-dimensional domain.

\subsubsection{Execution Handoff}
\label{sec:res-ver-tst-interface-location-handoff}

The Software workstream shall permit interface locations to be changed
without modifying either governing solver. Mesh generation, mapping,
checkpoint extraction, and test execution shall be parameterised by the
interface definition.

The verification software shall execute all candidate interface
locations, wait until every required branch has completed, and generate
a comparative report containing field, flux, spectral, and engineering
metrics.

\subsection{Governing Relations and Quantitative Definitions}
\label{sec:res-ver-tst-governing-relations}

% TODO: Record relevant equations, dimensionless groups, empirical models,
% extraction rules or quantitative definitions.
%
% Use align environments for displayed equations.
% Every displayed equation must have a unique \label{eq:...}.
% Do not place punctuation after displayed equations.

\subsection{Test Objective}
\label{sec:res-ver-tst-test-objective}

% TODO: Complete this RES-VER Domain-specific subsection.

\subsection{Reference or Benchmark Solution}
\label{sec:res-ver-tst-reference-benchmark-solution}

% TODO: Complete this RES-VER Domain-specific subsection.

\subsection{Test Configuration}
\label{sec:res-ver-tst-test-configuration}

% TODO: Complete this RES-VER Domain-specific subsection.

\subsection{Measured Quantities}
\label{sec:res-ver-tst-measured-quantities}

% TODO: Complete this RES-VER Domain-specific subsection.

\subsection{Error Metrics}
\label{sec:res-ver-tst-error-metrics}

% TODO: Complete this RES-VER Domain-specific subsection.

\subsection{Tolerance and Acceptance Criteria}
\label{sec:res-ver-tst-tolerance-acceptance-criteria}

% TODO: Complete this RES-VER Domain-specific subsection.

\subsection{Execution Handoff}
\label{sec:res-ver-tst-execution-handoff}

% TODO: Complete this RES-VER Domain-specific subsection.

\subsection{Candidate Approaches}
\label{sec:res-ver-tst-candidate-approaches}

% TODO: Define the credible candidate approaches considered.

\subsection{Critical Comparison}
\label{sec:res-ver-tst-critical-comparison}

% TODO: Compare candidates using physically and project-relevant criteria.
% Do not select an approach solely on popularity or implementation ease.

\subsection{Assumptions and Applicability}
\label{sec:res-ver-tst-assumptions}

% TODO: State assumptions and the conditions under which they are valid.

\subsection{Limitations and Failure Conditions}
\label{sec:res-ver-tst-limitations}

% TODO: State where the evidence, model or selected approach ceases to be
% reliable or applicable.

\subsection{Project-Specific Conclusions}
\label{sec:res-ver-tst-conclusions}

% TODO: Convert the evidence into conclusions specific to the Studentship
% project. Do not merely restate literature findings.

\subsection{Selected Approach and Rationale}
\label{sec:res-ver-tst-selected-approach}

% TODO: Record the selected approach only after sufficient evidence exists.
% State the alternatives rejected and the reasons for rejection.

\subsection{Unresolved Decisions}
\label{sec:res-ver-tst-unresolved-decisions}

% TODO: Record unresolved questions, required evidence and decision owners.

\subsection{Requirements and Handoffs}
\label{sec:res-ver-tst-requirements}

% TODO: State requirements passed to other Research Domains, Software,
% verification activities or later execution work.

\subsection{Deliverable Completion Record}
\label{sec:res-ver-tst-completion-record}

% TODO: Record whether the Deliverable is:
%
% - Not started
% - In progress
% - Draft complete
% - Reviewed
% - Accepted
%
% Acceptance requires:
%
% - the research questions to be answered;
% - claims to be supported by appropriate evidence;
% - assumptions and limitations to be explicit;
% - the project-specific conclusion to be stated;
% - downstream requirements to be traceable;
% - unresolved decisions to be closed or formally deferred.
